\section{UPPERCASE, lowercase, CamelCase}\label{sec:CaseOfNames}
Use \textit{CamelCase}, starting with a capital letter, for the identifiers of all kinds of classes: 
\begin{itemize}[nosep]
	\item{\lstinline|String|}
	\item{\lstinline|LocalDateTime|}
	\item{\lstinline|AccountOwner|}
	\item{…}
\end{itemize}

For the identifiers of local variables, methods and formal parameters, use \textit{camelCase} starting with a lowercase character: 
\begin{itemize}[nosep]
	\item{\lstinline|i|}
	\item{\lstinline|proceed|}
	\item{\lstinline|isValid|}
	\item{\lstinline|retValue|}
	\item{\lstinline|result|}
	\item{\lstinline|lastOwner|}
	\item{\lstinline|indexOfLastRecord|}
	\item{…}
\end{itemize}

Attributes are a bit special\footnote{see chapter \tqfullref{sec:AttributeNames}}: their identifiers are prefixed with “\verb#m_#”, followed by the concrete name in \textit{CamelCase} (with a leading capital letter):
\begin{itemize}[nosep]
	\item{\lstinline|m_Name|}
	\item{\lstinline|m_IsValid|}
	\item{\lstinline|m_LastOwner|}
	\item{\lstinline|m_IndexOfLastRecord|}
	\item{…}
\end{itemize}

For the identifiers of packages and modules, all \textit{lowercase} is preferred, although this is not mandatory:
\begin{itemize}[nosep]
	\item{\lstinline|java.lang|}
	\item{\lstinline|org.tquadrat.foundation|}
	\item{\lstinline|com.ibm|}
	\item{…}
\end{itemize}
In fact, these are 7~package \textit{names} (\lstinline|java|, \lstinline|lang|, \lstinline|org|, \lstinline|tquadrat|, \lstinline|foundation|, \lstinline|com| and \lstinline|ibm|) although we probably have only 3~\textit{packages}.

Nevertheless, you should avoid a package name starting with a capital letter, as this could be mixed up easily with a class name.

The \ccjpl demands to use all \textit{UPPERCASE} for constants and separating words with underscores, and usually this is implicitly expanded to \lstinline|enum| values, too\footnote{The \ccjpl ~were last updated in 1999, \lstinline|enum|s had been introduced with Java~5, and that was released at 2004-09-29~…}. I prefer a more relaxed approach for that – see chapters \tqfullvref{sec:Constants} and \tqfullvref{sec:EnumValues}.

\keeplisting{5}
Samples for constants are:
\begin{itemize}[nosep]
	\item{\lstinline|EOF|}
	\item{\lstinline|EMPTY_STRING|}
	\item{\lstinline|XML_ATTRIBUTE_LastAccessed|}
	\item{…}
\end{itemize}

\keeplisting{5}
Examples for an \lstinline|enum| named \lstinline|LightStatus| could be
\begin{itemize}[nosep]
	\item{\lstinline|LIGHT_STATUS_Unknown|}
	\item{\lstinline|LIGHT_STATUS_Off|}
	\item{\lstinline|LIGHT_STATUS_On|}
	\item{…}
\end{itemize}

Details are discussed in the chapters about the respective code elements, below.

%==============================================================================
% End of file!!
%==============================================================================
